\chapter*{Preface}
\addcontentsline{toc}{chapter}{Preface}
\markboth{Preface}{Preface}

A database engine is a small number of ideas, applied with care. Data lives in
fixed-size pages because disks move blocks. A buffer pool keeps the hot pages
in memory because memory is fast and small. A B+ tree turns a key into a page
in three or four reads because its fanout is in the hundreds. A lock manager
makes concurrent transactions look sequential. A write-ahead log, and an
algorithm called ARIES, make a committed transaction survive a crash that
happens at any instant.

This book explains those ideas from the ground up, and it explains each one
twice: first as theory, then as it appears in a real system. The real system
is \textbf{OntoDB}, a database management system for ontologies. Its data is
RDF, a graph of subject--predicate--object triples, and its query language is
SPARQL. Every chapter ends its theory with an \emph{In OntoDB} box that shows
what OntoDB does, and design-decision boxes that show what it chose instead,
and why.

\section*{Why an ontology database}

A relational engine teaches the same ideas, so why a triple store? Two reasons.
First, a triple store strips the problem to its core: there is one table, with
three columns, and every interesting question is a join. That makes indexing
and join ordering impossible to wave away. Second, ontologies are where
correctness matters most and is checked least: a scientific ontology edited by
many hands needs transactions, a log, and recovery just as much as a bank
ledger does.

\section*{How to read this book}

The parts follow the path of a request through the engine, from the bottom up.

\begin{description}[leftmargin=1.5em]
\item[Foundations] What an engine is made of, and enough RDF and SPARQL to follow the examples.
\item[Storage] Pages, the buffer pool, B+ trees, how RDF is indexed, and the log-structured alternative.
\item[Query processing] How a SPARQL query becomes a tree of operators, and how joins are chosen.
\item[Transactions] What isolation means, and how locking provides it.
\item[Recovery] Logging, the write-ahead rule, ARIES, and how OntoDB makes its databases durable today.
\end{description}

\section*{The boxes}

Four kinds of box recur. A \emph{Key idea} is the one thing to keep from a
section. \emph{In OntoDB} shows how the running example applies the theory. A
\emph{Design decision} names a choice OntoDB made, the alternatives, why it chose
as it did, and what the choice costs; a list of them follows the contents, and
they are the best place to start if you want to know why the engine is built the
way it is. \emph{Status} says what runs today. Each chapter ends with questions;
the answers are in \cref{app:answers}.

Code appears only where a picture or a sentence would be less precise: a key
layout, a log record, a few lines of pseudocode for an algorithm. The aim is to
understand \emph{why}, not to read an implementation.

\begin{status}
OntoDB is built feature by feature, and parts of this book describe designs
that are specified but not yet implemented. Each chapter's \emph{Status} box
says which. At the time of writing, the SPARQL shell, the triple key, the disk manager, the log
records, the log manager, and durable databases (snapshots, a checksummed write-ahead log,
recovery, and verification) run; the buffer pool, the B+ tree, the indexed
store, the optimizer, concurrency control, and page-level ARIES are designed
and being built.
\end{status}

\section*{Acknowledgments}

The structure of this book owes much to Alex Petrov's \emph{Database Internals},
whose division into storage engines and distributed systems it follows for the
first half. The recovery chapters follow Mohan et al.'s ARIES paper closely,
because nothing explains ARIES better than ARIES.
